% !TeX program = lualatex
% Compile twice: lualatex 227.tex
% All references and prose are embedded; no BibTeX database or external inputs.
\documentclass[11pt,a4paper]{article}
\usepackage[margin=24mm,headheight=14pt]{geometry}
\usepackage{fontspec}
\setmainfont{Latin Modern Roman}
\setsansfont{Latin Modern Sans}
\setmonofont{Latin Modern Mono}
\usepackage{amsmath,amssymb,mathtools}
\usepackage{unicode-math}
\setmathfont{Latin Modern Math}
\usepackage{microtype}
\usepackage{enumitem,xurl,needspace}
\usepackage[dvipsnames]{xcolor}
\usepackage{hyperref}
\hypersetup{unicode=true,colorlinks=true,linkcolor=MidnightBlue,urlcolor=MidnightBlue,
 pdftitle={Research Notebook: Section 227},pdfauthor={Research notebook},bookmarksnumbered=true}
\usepackage{bookmark}
\usepackage{tocloft}
\setlength{\cftsecnumwidth}{3em}
\cftsetpnumwidth{2.5em}
\cftsetrmarg{4em}
\usepackage{titlesec}
\titleformat{\section}{\Large\bfseries\raggedright}{\thesection}{0.7em}{}
\usepackage{fancyhdr}
\pagestyle{fancy}\fancyhf{}
\fancyhead[L]{\small\sffamily Open Problems | Research notebook}
\fancyhead[R]{\small 227 | 27 September 2026}
\fancyfoot[C]{\thepage}
\setlength{\parindent}{0pt}
\setlength{\parskip}{5pt plus 1pt}
\setlength{\emergencystretch}{3em}
\setcounter{tocdepth}{1}
\setcounter{secnumdepth}{1}
\setlist[itemize]{leftmargin=3em,itemsep=5pt,topsep=4pt,parsep=0pt}
\allowdisplaybreaks
\newcommand{\N}{\mathbb{N}}
\newcommand{\Z}{\mathbb{Z}}
\newcommand{\Q}{\mathbb{Q}}
\newcommand{\R}{\mathbb{R}}
\newcommand{\C}{\mathbb{C}}
\newcommand{\F}{\mathbb{F}}
\newcommand{\A}{\mathbb{A}}
\newcommand{\PP}{\mathbb P}
\newcommand{\E}{\mathbb{E}}
\newcommand{\eps}{\varepsilon}
\newcommand{\dd}{\,\mathrm d}
\newcommand{\sref}[2]{\hyperlink{source:#1:#2}{[S#2]}}
\newcommand{\eref}[2]{\hyperlink{extra:#1:#2}{[E#2]}}
% Turn the next switch off for a shorter reading copy. The quotations preserve
% every exactTarget field, including qualifications omitted from a short summary.
\newif\ifshowcatalogue
\showcataloguetrue
\newcommand{\cataloguescope}[1]{\ifshowcatalogue
 \paragraph{Exact scope in the supplied catalogue.}
 \begin{quote}\small #1\end{quote}\fi}

% Independently compilable extraction; original section text retained.
\numberwithin{equation}{section}
\begin{document}
\setcounter{section}{226}
% ================================================================
% problemId: problem.deterministic-parallel-perfect-matching-in-nc
% source releaseRank: 227; source releaseStatus: open_with_solved_subcases
% exactTarget SHA-256: 288b7b6cebbed74d641638fbfaf0c9f127f7e57ebdd79adc40e023286a206cd2
\clearpage
\section{\texorpdfstring{Perfect-Matching Search in Deterministic NC (General Graphs)}{Perfect-Matching Search in Deterministic NC (General Graphs)}}\label{227}
\subsection{Definitions and mathematical statement}
A perfect matching in a finite simple undirected graph is a set of disjoint edges covering every vertex. On the adjacency-matrix encoding of a labeled $n$-vertex graph, seek one logspace-uniform family of deterministic Boolean circuits that outputs a matching matrix, or a distinguished NO flag precisely when no perfect matching exists. With fan-in at most two for AND and OR, the bounds are
\[
 \operatorname{size}(C_n)\le Cn^c,\qquad
 \operatorname{depth}(C_n)\le C(\log(n+2))^k
\]
for constants independent of $G,n$. A logarithmic-space transducer must print $C_n$ from $1^n$, including its designated outputs. This is a multi-output search problem on general graphs, not just testing existence or treating bipartite graphs.
\par\smallskip\textit{Formulation and concept references:} \sref{227}{1}.
\cataloguescope{Assert the existence of one logspace-uniform family (C\_n) of deterministic Boolean AND/OR/NOT circuits which, for every n>=1 and every finite simple undirected graph G on labelled vertices 1,...,n given by its n-by-n adjacency matrix, outputs a perfect matching of G if one exists and outputs a distinguished NO flag if none exists. Matching output is an n-by-n edge-indicator matrix, required to encode disjoint edges covering every vertex and belonging to G; its bits are ignored when NO is output. Circuit size is bounded by C n\textasciicircum{}c and depth by C (log(n+2))\textasciicircum{}k for constants C,c,k independent of n and G, and AND/OR fan-in is at most two. A deterministic logspace transducer on unary 1\textasciicircum{}n outputs the complete circuit description, including all designated outputs. This fixes ordinary logspace-uniform NC for the multi-output search relation. It does not conflate a search algorithm with an unweighted decision algorithm or restrict the input to bipartite or planar graphs.}
\subsection{Short English statement}
Can perfect matchings be found deterministically in general graphs by uniform polynomial-size circuits of polylogarithmic depth?
% BEGIN RESEARCH ATTEMPT 227
\subsection{Research attempt}

\paragraph{Current frontier.}
Svensson--Tarnawski obtain deterministic quasi-NC for general-graph
matching: polylogarithmic parallel time with quasipolynomially many
processors. Polynomial size is the missing requirement, not merely
polylogarithmic depth. Their use of odd-set constraints and laminar faces
cannot be replaced by a bipartite argument without further work.
\sref{227}{1}, Abstract and Introduction.
The isolating-weights method originates in Mulmuley--Vazirani--Vazirani's
randomized parallel algorithm. \eref{227}{1}.
A 2025 manuscript revised in August 2026 gives a randomized
matrix-multiplication-time algorithm for the red-edge-count version of
\emph{Exact Matching}; it is a different problem and does not supply
the deterministic uniform circuits here. \eref{227}{3}, Abstract.
The attempt below isolates the uniform construction step instead of
claiming that a probabilistically existing list is already an algorithm.

\paragraph{Attack route.}
A decision-to-search self-reduction can require linearly many dependent
rounds. A parallel algebraic extraction is possible once a common
minimum matching is isolated; without isolation different processors
can make incompatible choices. The route pursued here combines a short
list of weight assignments, parallel determinant computations and an
explicit verification circuit. The main question becomes whether the
list can be generated uniformly, with both its length and its weight
magnitudes polynomially bounded. We first prove its nonuniform existence
and then show exactly why that does not close the target.

\paragraph{Attempt.}
\textbf{1. State the missing list construction precisely.}
Consider the following sufficient hypothesis, for some fixed $c,k$:
there is a logspace-uniform deterministic circuit family of polynomial
size and depth $O(\log^k n)$ which, from the adjacency matrix of $G$,
outputs a list of at most $n^c$ assignments
\[
 w^{(r)}:E(G)\longrightarrow\{1,\ldots,n^c\}.
\]
Whenever $G$ has a perfect matching, at least one assignment makes its
minimum-weight perfect matching unique. No requirement is imposed on a
list entry that fails to isolate. Call this hypothesis (ULI).
A stronger, input-independent construction would print such a list on
the complete graph for each $n$, isolating every subgraph with a perfect
matching. That stronger version also suffices, but is not assumed
necessary for deterministic NC.

\textbf{2. Isolated minimum matching gives parallel search, not just decision.}
Fix one assignment $w$, and form the skew-symmetric polynomial Tutte
matrix
\[
 T_w(t)_{ij}=\begin{cases}
 t^{w_{ij}},&i<j,\ ij\in E(G),\\
 -t^{w_{ij}},&i>j,\ ij\in E(G),\\
 0,&\text{otherwise}.
 \end{cases}
\]
For even $n$, its Pfaffian is
\[
 \operatorname{Pf}(T_w(t))=
    \sum_{M\text{ perfect}}\operatorname{sgn}(M)t^{w(M)},
 \qquad \det T_w(t)=\operatorname{Pf}(T_w(t))^2.
\]
If $M_*$ is the unique minimum with weight $W$, the coefficient of
$t^W$ in the Pfaffian is $1$ or $-1$, so
\begin{equation}\label{227:valuation}
 \operatorname{val}_t\det T_w=2W.
\end{equation}
Give the zero polynomial valuation $+\infty$. For every edge $e$,
compute also the determinant after deleting $e$. Then
\begin{equation}\label{227:extract}
 e\in M_*
 \quad\Longleftrightarrow\quad
 \operatorname{val}_t\det T_{w,G-e}>\operatorname{val}_t\det T_{w,G}.
\end{equation}
If $e\notin M_*$, the unique least term survives. If $e\in M_*$,
all remaining matchings have strictly greater weight, or none exist.
Cancellation among these higher terms cannot create a smaller exponent.
This proves \eqref{227:extract}, including the zero-polynomial case.
For a nonisolating assignment the extracted edge set may be wrong, so
we do not accept it without checking it.

All degrees here are at most $D=n\max w$, a polynomial. One concrete
way to perform the polynomial determinant computations is to evaluate
at $t=0,1,\ldots,D$ in parallel and interpolate. Each matrix entry has
polynomial bit length, and each determinant has bit length at most
$O(n\max w\log(D+1)+n\log n)$. Uniform parallel determinant algorithms
are available for these integer matrices. \eref{227}{2}, Section~2,
Definition~2.6 and Lemma~2.7. For interpolation one may use
\[
 P(t)=\sum_{j=0}^{D}P(j)
       \frac{\prod_{0\le r\le D,\ r\ne j}(t-r)}
            {(-1)^{D-j}j!(D-j)!}.
\]
A common denominator $D!$ makes all computations integral until a final
exact division; coefficient lists, factorials and products have
polynomial length and bit size. Balanced addition/multiplication trees
and standard uniform integer-arithmetic circuits have polylogarithmic
depth. Thus no arithmetic operation on an exponentially long integer is
being counted as one Boolean gate.

For every list entry compute \eqref{227:extract} in parallel, reject it
if its original determinant is zero, and verify that the resulting
indicator matrix uses graph edges and has exactly one incident selected
edge at every vertex. Those tests are uniform NC. Select the first
valid entry using prefix OR circuits; if none is valid, output NO.
Under (ULI), an existing perfect matching produces at least one valid
entry. If no perfect matching exists, no indicator matrix can pass the
verification. For odd $n$, output NO directly. This proves that (ULI)
implies the precise multi-output, logspace-uniform NC target. All
large fan-in conjunctions used in this description are replaced by
balanced fan-in-two trees.

\textbf{3. A polynomial universal list exists nonuniformly.}
Let $m=\binom n2$, $n\ge2$, and independently assign each of the $m$
possible edges a weight uniformly in $\{1,\ldots,2m\}$.
For any fixed nonempty family $\mathcal F$ of edge subsets, its
minimum-weight member is unique with probability at least $1/2$.
Here is the relevant isolation proof. Fix all weights but that of edge
$e$. Let $a$ be the least sum excluding $e$ among members containing
$e$, and $b$ the least weight of a member not containing $e$; if either
class is empty it cannot cause a tie. There is at most one critical
value $w_e=b-a$. If two different global minimizers exist, an edge in
their symmetric difference must have this critical value. A union bound
over the $m$ possible edges gives failure probability at most
$m/(2m)=1/2$. This is the isolating lemma of \eref{227}{1}, specialized
with all constants exposed.

Take $m+1$ independent assignments. For any fixed graph with a perfect
matching, the probability that every assignment fails to isolate is
at most $2^{-(m+1)}$. There are at most $2^m$ labelled graphs. Hence
\[
 \Pr(\text{some graph is not isolated by the entire list})
       \le2^m2^{-(m+1)}<1.
\]
There therefore exists, for each $n$, a universal list of $m+1$
assignments with weights at most $2m$. Its description length is
$O(n^4\log n)$. Hardwiring this list into the circuits of part~2 gives
\emph{nonuniform} polynomial-size, polylogarithmic-depth search circuits.
It does not give a logspace transducer printing their descriptions.
Enumerating all lists and testing all graphs would eventually find
one, but its time, space and circuit size are not those in the target.
The attempt stops at that uniformity gap, rather than promoting the
union bound to a deterministic parallel construction.

\textbf{4. Two ways to break an apparent shortcut.}
First, weights $w_e=2^{\operatorname{index}(e)}$ distinguish every
subset and therefore isolate every nonempty family. Their binary
encodings have only $O(n^2)$ bits, but the degrees of $T_w(t)$ are
exponential. The coefficient-based computation above would have
exponentially many output positions. A short encoding of an exponent
is not a polynomial degree bound. A compressed-polynomial valuation
algorithm with the same NC bounds is not supplied here.

Second, on $K_4$ give every edge weight one. With vertices in their
usual order,
\[
 \operatorname{Pf}(T_w)=t^2-t^2+t^2=t^2.
\]
Deleting edge $12$, $34$, $14$ or $23$ makes the Pfaffian zero; deleting
$13$ or $24$ leaves $2t^2$. The valuation rule would consequently
select the four-cycle $\{12,23,34,14\}$, not a matching. The companion
check verifies all six deletions exactly. This refutes the rule
\emph{without its isolation hypothesis}; it also explains the need for
a final output validator even when a list is guaranteed to contain one
good entry.

\paragraph{Stress test.}
The Pfaffian is used over characteristic zero, so its signs are not
silently discarded. Its polynomial identity does not assert that a
single arbitrary numeric substitution witnesses every perfect matching.
The determinant routine's ring-operation bounds have been accompanied
by degree and bit bounds. All deletions are computed in parallel, not
adaptively after an earlier choice of an edge.

The probabilistic argument proves existence of short advice, not uniform
NC. Conversely, (ULI) is sufficient but has not been proved necessary:
a future deterministic matching algorithm might avoid isolation entirely.
Nothing here makes an unweighted decision procedure automatically a
polylogarithmic-depth search procedure. Nor has the quasipolynomial
resource bound in \sref{227}{1} been relabelled polynomial.

\paragraph{Outcome.}
\textbf{Outcome: Conditional reduction.}

A uniform short-list isolation hypothesis has been shown to imply the
exact search relation in the catalogue, including NO outputs, Boolean
bit costs and logspace uniformity. The same method yields an elementary
nonuniform existence proof with an explicit polynomial advice length.
An exact four-vertex test rules out dropping isolation. These are
reconstructions and deductions within the notebook, not claims of a new
matching complexity theorem.

The missing step is a logspace-uniform polynomial-resource construction
of the lists, or a different search method avoiding them. The existence
proof by random choice does not provide that construction.
% END RESEARCH ATTEMPT 227
\subsection{Sources}
\begin{itemize}\raggedright
\item[\textnormal{[S1]}]\hypertarget{source:227:1}{} Svensson and Tarnawski, The Matching Problem in General Graphs is in Quasi-NC, introduction; FGT Section 2.1 defines Search-PM; Nanongkai–Scquizzato Section 2.2 fixes logspace-uniform NC.
\url{https://arxiv.org/pdf/1704.01929}.
{\small\textit{Catalogue formulation source.}}
% BEGIN ADDED REFERENCES 227
\item[\textnormal{[E1]}]\hypertarget{extra:227:1}{}
K. Mulmuley, U. V. Vazirani and V. V. Vazirani,
\emph{Matching is as easy as matrix inversion}, Combinatorica
\textbf{7} (1987), 105--113; isolating lemma and randomized parallel
matching method. \url{https://doi.org/10.1007/BF02579206}.
\item[\textnormal{[E2]}]\hypertarget{extra:227:2}{}
M. Soltys, \emph{Berkowitz's algorithm and clow sequences}, Electronic
Journal of Linear Algebra \textbf{9} (2002), 42--54, Section~2,
Definition~2.6 and Lemma~2.7. \url{https://arxiv.org/pdf/math/0201315}.
Used for the uniform, division-free parallel determinant construction.
\item[\textnormal{[E3]}]\hypertarget{extra:227:3}{}
R. Sato and Y. Yamaguchi, \emph{Exact Matching in Matrix Multiplication
Time}, arXiv:2508.04081v2, 10 August 2026, Abstract.
\url{https://arxiv.org/abs/2508.04081}.
The red-edge-count problem and randomized sequential time in this source
are not the uniform deterministic NC target. All consulted
27 September 2026.
% END ADDED REFERENCES 227
\end{itemize}

\end{document}
